% Created 2026-04-01 Wed 23:08
% Intended LaTeX compiler: pdflatex
\documentclass[11pt, letterpaper]{article}

\usepackage[utf8]{inputenc}
\usepackage[T1]{fontenc}
\usepackage{graphicx}
\usepackage{longtable}
\usepackage{wrapfig}
\usepackage{rotating}
\usepackage[normalem]{ulem}
\usepackage{amsmath}
\usepackage{amssymb}
\usepackage{capt-of}
\usepackage{hyperref}
\usepackage[margin=2.5cm]{geometry}      % Márgenes decentes
\usepackage[utf8]{inputenc}
\usepackage{palatino}                   % Tipografía elegante
\usepackage{xcolor}                     % Colores personalizados
\author{Equipo Alpine White}
\date{\textit{{[}2026-03-14 Sat]}}
\title{Bitácora Semanal 6: Administración de Sistemas Unix}
\hypersetup{
 pdfauthor={Equipo Alpine White},
 pdftitle={Bitácora Semanal 6: Administración de Sistemas Unix},
 pdfkeywords={},
 pdfsubject={},
 pdfcreator={},
 pdflang={English}}
\begin{document}

\maketitle
\setcounter{tocdepth}{2}
\tableofcontents

\section{Información General}
\label{sec:orga430dcd}
\begin{itemize}
\item \textbf{Semana:} Del 9 de Marzo 2026 al 13 de Marzo 2026
\item \textbf{Equipo:}
\begin{itemize}
\item Arreguín Salgado Gael Emiliano
\item Gramer Muñoz Omar Fernando
\item López Pérez Mariana
\item Nieto Gallegos Isaac Julián
\end{itemize}
\end{itemize}

En esta semana el enfoque se baso en crear paquetes para debian, algunos conceptos de red, ssh y maquinas virtuales
\subsection{Responsabilidades:}
\label{sec:orgf4205cb}

\begin{itemize}
\item Gael Emiliano: Notas físicas y aportaciones en nuevos conceptos y maquetado del reporte.
\item Omar Fernando: Notas físicas y aportaciones en contenido.
\item Mariana: Desarrollo de temas principales y estructura principal.
\item Isaac Julián: Cambios menores, primeros dos temas y algunas investigaciones en Secciones Libres
\end{itemize}
\section{Archivo /etc/hosts}
\label{sec:org39dbf4d}

El archivo \textbf{/etc/hosts} es un archivo de configuración local utilizado por el sistema
operativo para asociar \textbf{nombres de host} con \textbf{direcciones IP}.

Nos permite realizar una resolución de nombres sin necesidad de consultar un servidor DNS.
\subsection{Estructura del archivo}
\label{sec:orgbf3cd88}

Cada línea del archivo tiene la siguiente estructura:

\begin{verbatim}
IP_address   hostname   alias1 alias2 ...
\end{verbatim}

Ejemplo:

\begin{verbatim}
127.0.0.1   localhost
192.168.1.10 juanito 
192.168.1.33 lupita 
192.168.1.20 impresora impresora-local
\end{verbatim}
\subsection{Funcionamiento}
\label{sec:org4022d38}

Cuando una aplicación necesita resolver un nombre de host:

\begin{enumerate}
\item El sistema revisa primero \textbf{/etc/hosts}
\item Si no encuentra coincidencias, consulta el DNS
\item Devuelve la dirección IP correspondiente
\end{enumerate}

Esto depende del archivo /etc/nsswitch.conf donde se define el orden de resolución de nombres.
\subsection{Conceptos de red relacionados ( repaso )}
\label{sec:orgef04ed7}

\begin{itemize}
\item \textbf{Hostname}: nombre asignado a una máquina dentro de una red.
\item \textbf{Dirección IP}: identificador numérico único de un dispositivo en la red.
\item \textbf{Resolución de nombres}: proceso de convertir un nombre de host en una dirección IP.
\item \textbf{DNS (Domain Name System)}: sistema distribuido que traduce nombres de dominio a direcciones IP.
\end{itemize}
\section{Ping usando nombres definidos en hosts}
\label{sec:orgccbb46e}

Usamos lo anterior en clase para hacer ping entre compañeros de la clase, si un nombre está definido en \textbf{/etc/hosts}, 
es posible usarlo directamente con el comando \textbf{ping}.

Ejemplo:

\begin{verbatim}
ping juanito
\end{verbatim}

El sistema buscará ``juanito'' en \textbf{/etc/hosts}, obtendrá la IP asociada y enviará los paquetes ICMP a esa dirección.
\section{Directorio \textasciitilde{}/.ssh}
\label{sec:orgcdac668}

El directorio /.ssh contiene archivos relacionados con la autenticación segura usando \textbf{SSH (Secure Shell)}.
Este directorio normalmente tiene permisos restrictivos para proteger las claves.
\subsection{Archivos comunes}
\label{sec:orga166107}

\begin{itemize}
\item \textbf{id\textsubscript{rsa}} → clave privada
\item \textbf{id\textsubscript{rsa.pub}} → clave pública
\item \textbf{authorized\textsubscript{keys}} → lista de claves públicas autorizadas para acceder a la cuenta
\item \textbf{known\textsubscript{hosts}} → lista de hosts remotos conocidos y sus claves
\item \textbf{config} → configuración personalizada para conexiones SSH
\end{itemize}
\subsection{Función}
\label{sec:org2cc610b}

Permite autenticación segura entre máquinas sin necesidad de contraseñas,
mediante criptografía de clave pública.
\section{ssh-keygen}
\label{sec:org853a0c3}

El comando:

\begin{verbatim}
ssh-keygen
\end{verbatim}

se utiliza para generar un par de claves criptográficas para autenticación SSH.
\subsection{Proceso}
\label{sec:org6cabc20}

Al ejecutar el comando, el programa solicita:

\begin{enumerate}
\item Ubicación donde guardar la clave
\end{enumerate}

\begin{verbatim}
Enter file in which to save the key (/home/user/.ssh/id_rsa)
\end{verbatim}

\begin{enumerate}
\item Una \textbf{passphrase} (frase de contraseña) opcional para proteger la clave privada.
\end{enumerate}

\begin{verbatim}
Enter passphrase
\end{verbatim}
\subsection{Resultado}
\label{sec:org6017813}

Se generan dos archivos:

\begin{itemize}
\item clave privada → id\textsubscript{rsa}
\item clave pública → id\textsubscript{rsa.pub}
\end{itemize}

La clave pública se puede copiar al servidor en:

\begin{verbatim}
~/.ssh/authorized_keys
\end{verbatim}

para permitir autenticación sin contraseña.
\subsection{Flujo de autenticación SSH}
\label{sec:org99eac07}

El proceso de autenticación mediante llaves SSH sigue un modelo de criptografía de clave pública. Esto permite verificar la identidad del usuario sin necesidad de enviar contraseñas a través de la red.

El proceso general funciona de la siguiente manera:

\begin{verbatim}
Usuario (cliente)                Servidor

   ssh usuario@servidor
        |
        |-----> Solicitud de conexión SSH
        |
        |<----- Servidor envía su clave pública
        |
Cliente verifica si el servidor está en known_hosts
        |
Cliente usa su clave privada
para firmar la solicitud
        |
        |-----> Firma criptográfica
        |
Servidor verifica con la clave pública
en ~/.ssh/authorized_keys
        |
        |<----- Acceso permitido
\end{verbatim}

Gracias a este mecanismo, el servidor puede verificar la identidad del usuario sin que la clave privada sea transmitida a través de la red, lo cual mejora significativamente la seguridad del sistema.
\section{Creación de repositorios en Debian}
\label{sec:org1c92da0}

Esta semana durante las ayudantías vimos como crear repositorios en debian, 
en debian el software se distribuye mediante \textbf{paquetes .deb}.

Un repositorio es una colección organizada de paquetes accesible mediante herramientas como:

\begin{itemize}
\item apt
\item apt-get
\item aptitude
\end{itemize}
\subsection{Estructura básica de un repositorio}
\label{sec:org98393ab}

Un repositorio contiene:

\begin{itemize}
\item paquetes .deb
\item índices de paquetes
\item archivos de metadatos
\end{itemize}

Ejemplo simplificado:

\begin{verbatim}
repo/
 ├── pool/
 │   └── paquetes .deb
 ├── dists/
 │   └── stable/
 │       └── main/
 │           └── binary-amd64/
 │               └── Packages
\end{verbatim}
\section{Herramientas de empaquetado en Debian}
\label{sec:org6ec8e36}

\subsection{dh\textsubscript{make}}
\label{sec:orgdfb25f2}

Herramienta que crea la estructura inicial necesaria para empaquetar software en Debian.
Se usa normalmente dentro del directorio del código fuente.

\begin{verbatim}
dh_make
\end{verbatim}

Genera el directorio:

\begin{verbatim}
debian/
\end{verbatim}

que contiene archivos de configuración del paquete.
\subsection{Flujo del sistema de paquetes Debian}
\label{sec:org43bbbbc}

En los sistemas basados en Debian, el software se distribuye mediante paquetes \textbf{.deb}. Estos paquetes contienen los archivos necesarios para instalar una aplicación dentro del sistema.

El flujo general de instalación funciona de la siguiente manera:

\begin{verbatim}
Repositorio Debian
        |
        |  apt update
        v
Lista de paquetes disponibles
        |
        |  apt install paquete
        v
Descarga paquete .deb
        |
        v
Instalación con dpkg
        |
        v
Archivos copiados al sistema

/bin
/usr
/etc
/lib

        |
        v
Registro en base de datos

/var/lib/dpkg
\end{verbatim}

Esto permite que el sistema lleve un control completo de los programas instalados, facilitando su actualización, eliminación o mantenimiento.
\subsection{debuild}
\label{sec:org355bf25}

Compila el paquete Debian a partir del código fuente y los archivos en el directorio \textbf{debian/}.

\begin{verbatim}
debuild
\end{verbatim}

Resultado:

\begin{itemize}
\item archivo \textbf{.deb}
\item archivos de control y firma
\end{itemize}
\section{dpkg}
\label{sec:orgbd11714}

\textbf{dpkg} (Debian Package) es la herramienta base del sistema de gestión de paquetes
de Debian y de muchas distribuciones derivadas como Ubuntu.

Su función principal es \textbf{instalar, eliminar, configurar e inspeccionar paquetes
en formato *.deb}. A diferencia de herramientas como \textbf{apt}, \textbf{dpkg no resuelve
automáticamente dependencias}.

Las herramientas de alto nivel como \textbf{apt} utilizan \textbf{dpkg} internamente para
realizar la instalación real de los paquetes.
\subsection{Estructura de un paquete .deb}
\label{sec:org9d96cc6}

Un paquete \textbf{.deb} es un archivo comprimido que contiene tres componentes principales:

\begin{enumerate}
\item \textbf{debian-binary}  
Archivo que indica la versión del formato del paquete Debian.

\item \textbf{control.tar.gz} (o control.tar.xz)  
Contiene información del paquete:
\begin{itemize}
\item nombre
\item versión
\item dependencias
\item scripts de instalación
\end{itemize}

\item \textbf{data.tar.gz} (o data.tar.xz)  
Contiene los archivos reales que se instalarán en el sistema.
\end{enumerate}

Ejemplo para inspeccionar un paquete:

\begin{verbatim}
dpkg-deb -I paquete.deb
\end{verbatim}

Listar los archivos que contiene:

\begin{verbatim}
dpkg-deb -c paquete.deb
\end{verbatim}
\subsection{Base de datos de paquetes}
\label{sec:org117b6a7}

dpkg mantiene una base de datos local con información sobre los paquetes
instalados en el sistema.

Esta base de datos se encuentra en:

\begin{verbatim}
/var/lib/dpkg/
\end{verbatim}
\subsection{Comandos básicos}
\label{sec:org28b28c1}

\subsubsection{Instalar un paquete}
\label{sec:org3781aed}

\begin{verbatim}
sudo dpkg -i paquete.deb
\end{verbatim}

Si existen dependencias faltantes, dpkg mostrará un error y el paquete quedará
en estado \textbf{no configurado}.

Para corregir dependencias podemos usar:

\begin{verbatim}
sudo apt -f install
\end{verbatim}
\subsubsection{Eliminar un paquete}
\label{sec:org878f43a}

Eliminar el paquete pero conservar archivos de configuración:

\begin{verbatim}
sudo dpkg -r paquete
\end{verbatim}

Eliminar completamente el paquete y sus configuraciones:

\begin{verbatim}
sudo dpkg -P paquete
\end{verbatim}
\subsubsection{Listar paquetes instalados}
\label{sec:org41eb4ab}

\begin{verbatim}
dpkg -l
\end{verbatim}
\subsection{Comandos comunes de gestión de paquetes}
\label{sec:org92722ad}

\begin{center}
\begin{tabular}{ll}
Comando & Descripción\\
\hline
apt update & Actualiza la lista de paquetes disponibles\\
apt upgrade & Actualiza los paquetes instalados\\
apt install paquete & Instala un paquete\\
apt remove paquete & Elimina un paquete\\
dpkg -i archivo.deb & Instala un paquete manualmente\\
dpkg -l & \\
\end{tabular}
\end{center}
\subsection{Relación con APT}
\label{sec:org6dff278}

El sistema de paquetes de Debian funciona en dos niveles:

Nivel bajo:
\begin{itemize}
\item \textbf{dpkg} → instala y gestiona paquetes locales
\end{itemize}

Nivel alto:
\begin{itemize}
\item \textbf{apt}
\item \textbf{apt-get}
\item \textbf{aptitude}
\end{itemize}

Estas herramientas se encargan de:

\begin{itemize}
\item resolver dependencias
\item descargar paquetes desde repositorios
\item luego ejecutar \textbf{dpkg} para instalar.
\end{itemize}

Ejemplo de flujo típico:

\begin{enumerate}
\item apt descarga paquete
\item apt verifica dependencias
\item apt ejecuta dpkg para instalar
\end{enumerate}
\section{DKMS (Dynamic Kernel Module Support)}
\label{sec:org9ecd078}

Sistema que permite recompilar automáticamente módulos del kernel cuando se instala un nuevo kernel.
Esto es útil para controladores que no están incluidos directamente en el kernel.

Ejemplo típico:

\begin{itemize}
\item drivers de GPU
\item drivers de red
\item módulos de virtualización
\end{itemize}
\section{Compartir carpetas entre máquina virtual y anfitrión}
\label{sec:org32e8156}

En la clase del viernes vimos que muchas plataformas de virtualización permiten compartir directorios 
entre el sistema anfitrión y la máquina virtual.
\subsection{VirtualBox}
\label{sec:org2c182f6}

\begin{enumerate}
\item Instalar \textbf{Guest Additions} dentro de la VM.
\item Configurar carpeta compartida en la configuración de la máquina.
\item Montar la carpeta dentro de Linux.
\end{enumerate}

Ejemplo:

\begin{verbatim}
sudo mount -t vboxsf carpeta_compartida /mnt/compartida
\end{verbatim}
\subsection{VMware}
\label{sec:orgc94263b}

\begin{enumerate}
\item Instalar \textbf{VMware Tools}
\item Activar \textbf{Shared Folders}
\item Las carpetas aparecen generalmente en:
\end{enumerate}

\begin{verbatim}
/mnt/hgfs/
\end{verbatim}
\subsection{Uso}
\label{sec:org06f3ac2}

Las carpetas compartidas permiten:

\begin{itemize}
\item transferir archivos entre sistemas
\item compartir código
\item acceder a documentos del anfitrión desde la VM
\end{itemize}
\end{document}
